% !TEX root = ../tfg_etsiinf_NombreApellidos.tex
\chapter{Trabajo relacionado y Estado del Arte} \label{chp:state-of-the-art}

Este cap{\'i}tulo revisa el estado del arte relacionado con el control de
veh{\'i}culos a{\'e}reos no tripulados mediante aprendizaje por refuerzo y con
las infraestructuras software y de simulaci{\'o}n necesarias para su
desarrollo. La revisi{\'o}n se organiza en torno a cuatro ejes. En primer
lugar, se presentan algunos marcos software relevantes en rob{\'o}tica a{\'e}rea
y, en particular, la evoluci{\'o}n desde Aerostack hasta Aerostack2. A
continuaci{\'o}n, se analizan las principales l{\'i}neas de trabajo en control
de cuadric{\'o}pteros mediante aprendizaje por refuerzo, diferenciando entre
control a bajo nivel, seguimiento de trayectorias y vuelo agresivo. En tercer
lugar, se estudian los simuladores y entornos de entrenamiento m{\'a}s
utilizados en la literatura. Finalmente, se sintetizan las limitaciones
observadas y se posiciona el presente trabajo respecto a ellas.

%%%%%%%%%%%%%%%%%%%%%%%%%%%%%%%%%%%%%%%%%%%%%%%%%%%%%%%%%%%%%%%%%%%%%%%%%%%%%%%%
%%%%%%%%%%%%%%%%%%%%%%%%%%%%%%%%%%%%%%%%%%%%%%%%%%%%%%%%%%%%%%%%%%%%%%%%%%%%%%%%

\section{Marcos software para rob{\'o}tica a{\'e}rea}

El desarrollo de sistemas a{\'e}reos aut{\'o}nomos requiere combinar de forma
coherente control, estimaci{\'o}n, percepci{\'o}n, planificaci{\'o}n y
coordinaci{\'o}n entre m{\'u}ltiples componentes heterog{\'e}neos. Este problema
de integraci{\'o}n ha sido identificado desde hace a{\~n}os como uno de los
principales obst{\'a}culos para trasladar algoritmos aislados a sistemas
rob{\'o}ticos completos \cite{sanchez2017aerostack}. En ese contexto,
S{\'a}nchez-L{\'o}pez et al.\ propusieron Aerostack como un marco software
abierto, modular y multicapa orientado al desarrollo de sistemas a{\'e}reos
aut{\'o}nomos complejos \cite{sanchez2017aerostack}. Su aportaci{\'o}n m{\'a}s
relevante fue proporcionar una arquitectura reutilizable, espec{\'i}ficamente
dise{\~n}ada para rob{\'o}tica a{\'e}rea, capaz de integrar desde funciones
reactivas de control hasta capas deliberativas y sociales para misiones de
mayor complejidad.

La evoluci{\'o}n natural de esta l{\'i}nea es Aerostack2, redise{\~n}ado sobre
ROS 2 y orientado a sistemas a{\'e}reos multi-robot m{\'a}s vers{\'a}tiles y
portables \cite{fernandez2023aerostack2}. Frente a la versi{\'o}n anterior,
Aerostack2 enfatiza la independencia de plataforma, la arquitectura basada en
\emph{plugins} y la definici{\'o}n de misiones a trav{\'e}s de comportamientos,
lo que facilita reutilizar componentes y conectar diferentes simuladores o
plataformas reales. Para este TFM, Aerostack2 no es solo una herramienta de
soporte, sino el sustrato tecnol{\'o}gico sobre el que se construye el entorno
experimental. Por ello, resulta especialmente relevante en la revisi{\'o}n del
estado del arte, ya que buena parte de los trabajos publicados en aprendizaje
por refuerzo para drones se apoyan en simuladores o \emph{stacks} propios, pero
no en ecosistemas modulares integrados con ROS 2 de uso directo en rob{\'o}tica
a{\'e}rea.

\section{Aprendizaje por refuerzo para el control de UAV}

Las revisiones recientes muestran que el aprendizaje por refuerzo profundo se ha
aplicado a problemas de planificaci{\'o}n, navegaci{\'o}n y control de UAV, con
especial inter{\'e}s en escenarios donde la din{\'a}mica del veh{\'i}culo es no
lineal o donde el dise{\~n}o manual del controlador resulta costoso
\cite{azar2021drone_review}. Al mismo tiempo, la literatura tambi{\'e}n
subraya limitaciones persistentes: tiempos de entrenamiento elevados,
dificultades de transferencia del simulador al mundo real y fuerte dependencia
de la calidad del entorno de simulaci{\'o}n \cite{azar2021drone_review,
chan2024drone_simulators}.

Una parte significativa de los trabajos se ha centrado en control a bajo nivel,
es decir, en aprender pol{\'i}ticas que act{\'u}an cerca de los actuadores o de
las consignas de actitud. Pi et al.\ propusieron un controlador de bajo nivel
basado en aprendizaje por refuerzo capaz de generalizar a veh{\'i}culos
multirrotor con distintas configuraciones f{\'i}sicas, tratando de reducir la
necesidad de reentrenamiento cuando cambian los par{\'a}metros del veh{\'i}culo
\cite{pi2021lowlevel}. En una l{\'i}nea relacionada, Yoo et al.\ combinaron
controladores deterministas cl{\'a}sicos con una pol{\'i}tica aprendida,
mostrando que una aproximaci{\'o}n h{\'i}brida puede mejorar la convergencia y
el seguimiento de trayectorias mediante control en velocidad
\cite{yoo2021hybrid}. Estos trabajos son especialmente relevantes para la
extensi{\'o}n futura planteada en esta tesis, en la que se contempla encadenar
un controlador de alto nivel con otro de menor nivel que produzca tasas
angulares.

Otra l{\'i}nea muy activa es la del vuelo agresivo y el empuje del dron hacia
sus l{\'i}mites din{\'a}micos. Kaufmann et al.\ demostraron en
\emph{Deep Drone Acrobatics} que es posible entrenar en simulaci{\'o}n una
pol{\'i}tica para realizar maniobras acrob{\'a}ticas y transferirla al dron real
sin ajuste fino posterior \cite{kaufmann2020acrobatics}. M{\'a}s adelante, los
mismos autores analizaron comparativamente distintas abstracciones de control y
observaron que las pol{\'i}ticas que act{\'u}an sobre \emph{body-rates} y thrust
resultan m{\'a}s robustas en transferencia que aquellas que operan
directamente sobre los motores \cite{kaufmann2022benchmark}. Este punto es de
gran inter{\'e}s para el presente trabajo, porque ayuda a entender que el nivel
de abstracci{\'o}n elegido para la acci{\'o}n no es una decisi{\'o}n menor, sino
un elemento central del dise{\~n}o del problema de aprendizaje.

Los resultados m{\'a}s recientes en vuelo agresivo refuerzan esta idea.
Eschmann et al.\ mostraron que, con una arquitectura adecuada y un simulador muy
optimizado, es posible entrenar controladores de bajo nivel en tiempos
extraordinariamente reducidos y mantener capacidad de transferencia a plataformas
reales \cite{eschmann2024seconds}. En paralelo, Song et al.\ estudiaron de
manera sistem{\'a}tica la competici{\'o}n entre control {\'o}ptimo y aprendizaje
por refuerzo en un problema de carrera aut{\'o}noma con drones, concluyendo que
RL puede superar al control {\'o}ptimo cuando el objetivo se formula
directamente a nivel de tarea y se entrena con suficiente robustez frente a
incertidumbres del modelo \cite{song2023racing}. Este resultado es
particularmente pertinente dado el contexto de A2RL que motiva la tesis.

Sin embargo, la literatura revisada tambi{\'e}n revela una cierta dispersi{\'o}n
en el tipo de problema abordado. Muchos trabajos se centran en estabilizaci{\'o}n
de actitud, control a nivel de actuadores, vuelo acrob{\'a}tico o navegaci{\'o}n
end-to-end basada en percepci{\'o}n \cite{kaufmann2020acrobatics,
pi2021lowlevel, yoo2021hybrid}. Por el contrario, es menos frecuente encontrar
trabajos orientados a un controlador intermedio de alto nivel, cuya entrada sea
una referencia de posici{\'o}n global y cuya salida sean velocidades lineales y
velocidad angular en yaw, integrados adem{\'a}s en un \emph{stack} rob{\'o}tico
como Aerostack2. Esta observaci{\'o}n es una inferencia a partir de las obras
revisadas y ayuda a justificar el espacio de contribuci{\'o}n de este TFM.

\section{Simuladores y entornos de entrenamiento para RL en drones}

El entrenamiento mediante aprendizaje por refuerzo en rob{\'o}tica a{\'e}rea
depende fuertemente del simulador empleado. Un simulador adecuado debe ofrecer
fidelidad din{\'a}mica suficiente para capturar el comportamiento del veh{\'i}culo,
pero tambi{\'e}n un coste computacional reducido que permita generar gran
cantidad de experiencia de entrenamiento \cite{chan2024drone_simulators}. La
literatura reciente coincide en que este equilibrio entre realismo, velocidad y
facilidad de integraci{\'o}n es uno de los factores determinantes del {\'e}xito
de una soluci{\'o}n basada en RL \cite{azar2021drone_review,
chan2024drone_simulators}.

Entre los entornos m{\'a}s influyentes se encuentra Flightmare, un simulador de
cuadric{\'o}pteros dise{\~n}ado para desacoplar el renderizado y la f{\'i}sica,
lo que permite combinar flexibilidad con una velocidad de simulaci{\'o}n muy
elevada \cite{song2021flightmare}. Su inter{\'e}s para aprendizaje por refuerzo
reside en que puede simular centenares de drones en paralelo, lo cual lo
convierte en una referencia cuando el objetivo es maximizar el rendimiento de
muestreo. Desde otra perspectiva, Air Learning ampli{\'o} el uso de AirSim con
un entorno tipo Gym para navegaci{\'o}n aut{\'o}noma, introduciendo
aleatorizaci{\'o}n del dominio, \emph{hardware-in-the-loop} y m{\'e}tricas de
``calidad de vuelo'' orientadas al despliegue en plataformas con recursos
limitados \cite{krishnan2021airlearning}. Este trabajo es relevante porque
muestra que el valor de un simulador para RL no depende solo de la din{\'a}mica,
sino tambi{\'e}n de su capacidad para aproximar las restricciones computacionales
del sistema final.

En una direcci{\'o}n m{\'a}s ligera y abierta, Panerati et al.\ propusieron un
entorno Gym basado en PyBullet para cuadric{\'o}pteros individuales y
multiagente, expl{\'i}citamente orientado a facilitar la comparaci{\'o}n entre
enfoques de control cl{\'a}sico y de aprendizaje \cite{panerati2021pybullet}. Su
aportaci{\'o}n es especialmente interesante para trabajos acad{\'e}micos porque
reduce la fricci{\'o}n de entrada, ofrece portabilidad y proporciona ejemplos
tanto de control PID como de tareas de RL. Finalmente, \emph{Learning to Fly in
Seconds} plantea un simulador extremadamente optimizado para entrenamiento de
controladores de bajo nivel, subrayando que la eficiencia computacional puede
llegar a ser tan importante como la fidelidad din{\'a}mica cuando el objetivo es
acelerar el ciclo de investigaci{\'o}n \cite{eschmann2024seconds}.

Las revisiones de conjunto confirman esta diversidad de soluciones. Chan et
al.\ comparan simuladores orientados a RL desde la perspectiva de sus
prestaciones, interfaces y capacidades de integraci{\'o}n, y concluyen que no
existe una herramienta universalmente superior, sino distintas elecciones
dependiendo del problema, del tipo de aprendizaje y del nivel de realismo
requerido \cite{chan2024drone_simulators}. Desde la perspectiva de este TFM,
esta conclusi{\'o}n es importante: el simulador de multirrotores de Aerostack2
puede no ofrecer el fotorrealismo de otras plataformas, pero su ligereza y su
integraci{\'o}n natural con ROS 2 lo hacen especialmente atractivo para un
entorno vectorizado de entrenamiento orientado al control din{\'a}mico.

\section{S{\'i}ntesis cr{\'i}tica y posicionamiento del trabajo}

La revisi{\'o}n anterior permite extraer varias conclusiones. En primer lugar,
el aprendizaje por refuerzo ha demostrado ser una herramienta v{\'a}lida para el
control de drones en problemas de estabilizaci{\'o}n, seguimiento, navegaci{\'o}n
y vuelo agresivo \cite{kaufmann2020acrobatics, pi2021lowlevel,
yoo2021hybrid, song2023racing}. En segundo lugar, el rendimiento final depende
de forma decisiva del nivel de abstracci{\'o}n de la acci{\'o}n, del dise{\~n}o
de la recompensa y de la calidad del simulador empleado
\cite{kaufmann2022benchmark, chan2024drone_simulators}. En tercer lugar, la
literatura ha prestado especial atenci{\'o}n a simuladores propios o
entornos Gym independientes, pero menos a la integraci{\'o}n del aprendizaje por
refuerzo con \emph{frameworks} modulares de rob{\'o}tica a{\'e}rea basados en
ROS 2 \cite{fernandez2023aerostack2, song2021flightmare,
krishnan2021airlearning, panerati2021pybullet}.

A partir de las fuentes revisadas, puede inferirse que existe una oportunidad de
investigaci{\'o}n en la construcci{\'o}n de entornos de aprendizaje por refuerzo
ligeros, vectorizables y directamente integrados con un \emph{stack} rob{\'o}tico
operativo como Aerostack2. Esta necesidad es todav{\'i}a m{\'a}s clara cuando el
problema de control no se formula a nivel de actuadores, sino a un nivel
intermedio m{\'a}s pr{\'o}ximo a la planificaci{\'o}n y a la ejecuci{\'o}n de
misiones: recibir una referencia global de posici{\'o}n y generar comandos de
velocidad compatibles con la arquitectura del dron, imponiendo adem{\'a}s una
restricci{\'o}n funcional sobre la orientaci{\'o}n en yaw para alinear el dron
con su direcci{\'o}n de movimiento.

En este punto se sit{\'u}a el presente Trabajo Fin de M{\'a}ster. La
contribuci{\'o}n propuesta no consiste solo en entrenar un controlador, sino en
sentar la base experimental necesaria para hacerlo dentro de Aerostack2: validar
la interfaz entre simulador, ROS 2 y entorno de aprendizaje, preparar un
entorno vectorizable y establecer posteriormente una comparaci{\'o}n rigurosa
frente a un controlador PID de referencia. Desde esta perspectiva, el trabajo se
alinea con la literatura m{\'a}s reciente sobre control aprendido para
cuadric{\'o}pteros, pero aporta un enfoque diferencial en t{\'e}rminos de
integraci{\'o}n software, reproducibilidad y adecuaci{\'o}n a un problema de
control jer{\'a}rquico progresivo.
